\section{Matching bipartito: Hopcroft--Karp}
\whenuse{Emparejar el m\'aximo n\'umero de elementos entre dos conjuntos con
compatibilidades binarias y sin costos.}

Los v\'ertices izquierdos son $[0,n)$ y los derechos $[0,m)$.

Una arista $(l,r)$ indica que el elemento izquierdo $l$ puede emparejarse con
el derecho $r$; cada v\'ertice puede pertenecer a lo sumo a una pareja.
Hopcroft--Karp busca caminos aumentantes: alterna aristas fuera/dentro del
matching para reemplazar parejas y aumentar su tama\~no en uno. El BFS agrupa
los caminos aumentantes m\'as cortos por niveles y los DFS procesan varios en la
misma fase. \texttt{maximum\_matching()} devuelve la cantidad de parejas;
\texttt{match\_left[l]} y \texttt{match\_right[r]} contienen su pareja o $-1$.

\complexity{$O(E\sqrt V)$ tiempo y $O(V+E)$ memoria.}

\begin{lstlisting}[style=cpp]
struct HopcroftKarp {
    int left_size, right_size;
    vector<vector<int>> graph;
    vector<int> match_left, match_right, distance;

    HopcroftKarp(int left_count, int right_count)
        : left_size(left_count), right_size(right_count), graph(left_count),
          match_left(left_count, -1), match_right(right_count, -1),
          distance(left_count) {}

    void add_edge(int left, int right) { graph[left].push_back(right); }

    bool bfs() {
        queue<int> q;
        fill(distance.begin(), distance.end(), -1);
        for (int left = 0; left < left_size; ++left)
            if (match_left[left] == -1) distance[left] = 0, q.push(left);
        bool augmenting_path = false;
        while (!q.empty()) {
            int left = q.front(); q.pop();
            for (int right : graph[left]) {
                int next_left = match_right[right];
                if (next_left == -1) augmenting_path = true;
                else if (distance[next_left] == -1) {
                    distance[next_left] = distance[left] + 1;
                    q.push(next_left);
                }
            }
        }
        return augmenting_path;
    }

    bool dfs(int left) {
        for (int right : graph[left]) {
            int next_left = match_right[right];
            if (next_left == -1 ||
                (distance[next_left] == distance[left] + 1 && dfs(next_left))) {
                match_left[left] = right;
                match_right[right] = left;
                return true;
            }
        }
        distance[left] = -1;
        return false;
    }

    int maximum_matching() {
        int matching = 0;
        while (bfs())
            for (int left = 0; left < left_size; ++left)
                if (match_left[left] == -1 && dfs(left)) ++matching;
        return matching;
    }
};
\end{lstlisting}

Ejemplo: asignar el m\'aximo n\'umero de personas a trabajos compatibles.

\begin{lstlisting}[style=cpp]
HopcroftKarp matching(3,3);
matching.add_edge(0,0); matching.add_edge(0,1);
matching.add_edge(1,1); matching.add_edge(2,1);
int pairs=matching.maximum_matching();
for(int person=0;person<3;++person)
    cout<<person<<" -> "<<matching.match_left[person]<<'\n';
\end{lstlisting}

\section{Asignaci\'on de costo m\'inimo: H\'ungaro}
\whenuse{Asignar cada elemento de un lado a uno distinto del otro minimizando
un costo total en una matriz completa o casi completa.}

Matriz rectangular con $n\le m$. Devuelve costo y columna asignada a cada fila.
Admite costos negativos.

\texttt{cost[i][j]} es el costo de asignar la fila $i$ a la columna $j$; todas
las filas reciben columnas distintas. El algoritmo incorpora una fila por vez y
usa potenciales para conservar costos reducidos no negativos mientras encuentra
la reasignaci\'on m\'as barata. Devuelve \texttt{\{costo,assignment\}}, donde
\texttt{assignment[i]} es la columna elegida para la fila $i$.

\complexity{$O(n^2m)$ tiempo y $O(n+m)$ memoria.}

\begin{lstlisting}[style=cpp]
pair<ll, vector<int>> hungarian(const vector<vector<ll>> &cost) {
    int n = cost.size();
    if (n == 0) return {0, {}};
    int m = cost[0].size();
    assert(n <= m && m > 0);

    vector<ll> row_potential(n + 1), column_potential(m + 1);
    vector<int> row_for_column(m + 1), previous_column(m + 1);
    for (int row = 1; row <= n; ++row) {
        row_for_column[0] = row;
        int column0 = 0;
        vector<ll> minimum(m + 1, LINF);
        vector<char> used(m + 1, false);
        do {
            used[column0] = true;
            int row0 = row_for_column[column0], column1 = 0;
            ll delta = LINF;
            for (int column = 1; column <= m; ++column) if (!used[column]) {
                ll current = cost[row0 - 1][column - 1]
                           - row_potential[row0] - column_potential[column];
                if (current < minimum[column])
                    minimum[column] = current, previous_column[column] = column0;
                if (minimum[column] < delta)
                    delta = minimum[column], column1 = column;
            }
            for (int column = 0; column <= m; ++column) {
                if (used[column]) {
                    row_potential[row_for_column[column]] += delta;
                    column_potential[column] -= delta;
                } else minimum[column] -= delta;
            }
            column0 = column1;
        } while (row_for_column[column0] != 0);

        do {
            int column1 = previous_column[column0];
            row_for_column[column0] = row_for_column[column1];
            column0 = column1;
        } while (column0 != 0);
    }

    vector<int> assignment(n, -1);
    for (int column = 1; column <= m; ++column)
        if (row_for_column[column])
            assignment[row_for_column[column] - 1] = column - 1;
    return {-column_potential[0], assignment};
}
\end{lstlisting}

Ejemplo: asignar tres personas a trabajos distintos minimizando el costo.

\begin{lstlisting}[style=cpp]
vector<vector<ll>> cost={{4,1,3},{2,0,5},{3,2,2}};
auto [minimum_cost,job]=hungarian(cost);
cout<<minimum_cost<<'\n'; // 5
for(int person=0;person<3;++person)
    cout<<person<<" -> "<<job[person]<<'\n';
\end{lstlisting}

\subsection{Elecci\'on r\'apida}

\begin{itemize}
    \item Compatibilidad s\'i/no y maximizar parejas: \textbf{Hopcroft--Karp}.
    \item Matriz de costos y asignar todas las filas: \textbf{H\'ungaro}.
    \item Capacidades, cupos o cantidad de asignaciones variable con costos:
    \textbf{Min-Cost Max-Flow}.
\end{itemize}
